

\chapter{Objetivos}\label{cap:obj}

O objetivo desse trabalho é desenvolver o protótipo de um sistema de alinhamento automático
de antena helicoidal impressa em impressora 3D. Esse protótipo contará com um sistema de controle motorizado e uma integração com sistemas externos,
como aplicativo de rastreamento de satélites.


\section{Objetivos Especificos}\label{sec:esp}
\subsection{Modelagem e Prototipagem}\label{sec:mod}
A fim de obter êxito no rastreio do sinal enviado pelos satélites, .
Primeiramente, será necessária uma modelagem dos ângulos e os modelos para permitir o cálculo da cinemática
direta do mecanismo e permitir o apontamento do sistema.

\subsection{Sistema de Controle}\label{sec:ctr}
Após a modelagem do sistema, será necessário calcular a resposta do sistema a
estímulos e estabilizar o controle do mecanismo. O sistema de controle utilizado será de
malha aberta, utilizando motores de passo para realizar o controle do mecanismo do sistema.
O sistema de controle deve ser estável e responsivo para acompanhar a trajetória do veículo selecionado. 
Além disso, será utilizado um filtro $\alpha-\beta-\gamma$ para permitir um rastreio 

\subsection{Características Mecâncias}\label{sec:ssb}
O mecanismo deverá possuir duas características principais: velocidade e precisão.
O primeiro fator está relacionado com a resposta do sistema à tarefa. Uma passagem de satélites
de órbita baixa é marcada por momentos de aceleração e desaceleração. O protótipo deverá ser capaz de
efetuar o apontamento da antena em todas as direções de uma semiesfera e ser sensível o suficiente para ter a precisão de
apontamento necessária para a comunicação.

\subsection{Interface com Sistemas Externos}
Diversos aplicativos, programas e bancos de dados, disponíveis na internet, fornecem informações sobre órbitas e mudanças de posição
de satélites. O envio desses parâmetros acontece por uma comunicação TCP (Transmission Control Protocol), pelo qual o
sistema irá receber ângulos de azimute e elevação para o apontamento da antena e irá retornar os dados de sua posição atual. Dessa forma, o protótipo deverá ser capaz de
receber e enviar parâmetros para realizar a comunicação com aplicativos e programas.